% ================================================================
%  PREGUNTAS EMPÍRICAS — FINANZAS PÚBLICAS DEL PERÚ
%  Fuente de datos: MMM 2026-2029 (C-1, C-6, C-10) y BCRP (C-4)
%  Insertar dentro del \begin{enumerate} principal del problema set
% ================================================================

        \begin{enumerate}[label = (\alph*), start=3]   % empieza en (c) porque (a) y (b) ya están definidas

            % ─────────────────────────────────────────────────────────────
            \item \textbf{La aritmética del déficit.}
            El Cuadro~1 del MMM reporta tres indicadores de resultado fiscal
            para el Sector Público No Financiero (SPNF):
            el \emph{resultado primario}, los \emph{intereses} y el
            \emph{resultado económico}.

            \begin{itemize}
                \item Identifica los tres valores (en \% del PBI) para
                      los años 2024 y 2025.
                \item ¿Qué relación aritmética existe entre ellos?
                      Escríbela como una identidad y verifícala con
                      los datos.
                \item A partir de las proyecciones 2026--2029, ¿en qué
                      año el resultado primario deja de ser negativo?
                      ¿Qué significa económicamente que el resultado
                      primario sea positivo pero el resultado económico
                      siga siendo negativo?
            \end{itemize}

            % ─────────────────────────────────────────────────────────────
            \item \textbf{Convencional vs.\ estructural: ¿qué esconde
                  el ciclo?}
            Además del resultado económico ``convencional'', el MMM
            reporta un \emph{resultado económico estructural}.
            En 2024, el resultado convencional fue de
            $-3{,}46\%$ del PBI y el estructural de
            $-3{,}94\%$ del PBI potencial.\footnote{Para el cálculo
            del resultado estructural el MEF utiliza la metodología
            aprobada por la Resolución Ministerial N.\textordmasculine~024-2016-EF/15.}

            \begin{itemize}
                \item ¿Por qué ambos indicadores difieren?
                      ¿Qué componentes de los ingresos y del PBI
                      se ``corrigen'' al pasar del resultado convencional
                      al estructural?
                \item El resultado estructural es \emph{más negativo}
                      que el convencional en 2024 y 2025, pero ambos
                      convergen hacia 2029.
                      ¿Qué te dice eso sobre la posición cíclica de la
                      economía peruana en esos años y sobre la
                      sostenibilidad de la consolidación fiscal proyectada?
            \end{itemize}

            % ─────────────────────────────────────────────────────────────
            \item \textbf{¿Qué explica la trayectoria del déficit?}
            El Cuadro~4 del repositorio contiene la serie histórica del
            resultado económico del SPNF entre 2000 y 2025 (en \%
            del PBI), construida con datos del BCRP.

            \begin{itemize}
                \item Describe la trayectoria del resultado económico
                      entre 2019 y 2025. Identifica al menos dos
                      episodios de cambio brusco y señala a qué años
                      corresponden.
                \item Para el episodio de deterioro 2022--2024
                      (de $-1{,}68\%$ a $-3{,}46\%$ del PBI),
                      ¿qué combinación de factores por el lado de los
                      ingresos y por el lado del gasto podría explicarlo?
                      Apóyate en los Cuadros~6 y~10 del repositorio.\footnote{Cuadro~6:
                      Ingresos del gobierno general (\% del PBI).
                      Cuadro~10: Gasto no financiero del gobierno
                      general (\% del PBI).}
                \item El MMM proyecta una mejora sostenida del resultado
                      económico hasta $-0{,}99\%$ del PBI en 2029.
                      Considerando la historia reciente, ¿qué condiciones
                      tendrían que cumplirse para que esa trayectoria
                      sea alcanzable?
                      ¿Cuáles son los principales riesgos que la
                      podrían desviar?
            \end{itemize}

        \end{enumerate}
